\section{Problem and data}

\textbf{Files.} Four CSV files define the task.
\begin{itemize}
  \item \texttt{train-test.csv}: \NTrain{} loads from 1 January to 31 October 2025, with the posted rate.
  \item \texttt{validation.csv}: \NValidation{} loads from 1 November to 31 December 2025, without rates. These are
        the loads we must price.
  \item \texttt{validation-predictions-template.csv}: the load IDs, in the order the submission must follow.
  \item \texttt{december-chart-inputs.csv}: one load (Lexington to Fort Wayne, 360 miles, Dry Van, 32,000 lb)
        for every day of December.
\end{itemize}

\textbf{Columns.} Each load has a pickup and delivery city with coordinates, the road distance, the equipment type
(Dry Van, Reefer or Flatbed), the weight, the date, a \texttt{market\_index} and a \texttt{quote\_signal}.

\textbf{Target.} The posted rate in dollars. We model its logarithm, so errors are relative and a 2\% miss counts
the same on a \usd{300} and a \usd{5,000} load. The forecast is the exponential of the predicted logarithm, which
estimates the median rate.

\textbf{Horizon.} On 31 October we know every label so far and must price loads 1 to 61 days ahead.
The market index of each target day is known, because it is a column of the scoring file.

\textbf{Unseen cities.} \UnseenCities{} cities in the scoring file never appear in training (\UnseenCityList).
They touch \UnseenShare\% of the scoring loads. A model that memorises city names cannot price them, so the model
uses coordinates instead.

\textbf{Metric.} The assessment does not name a scoring metric. We chose the mean absolute error in dollars
(Section~\ref{sec:validation}) and report several others next to it.
